\chapter*{Abstract}
\addcontentsline{toc}{chapter}{Abstract}

Mutual information (MI) measures how much observing one variable reduces
uncertainty about another. This thesis considers how to test whether MI
between two categorical variables differs across independent populations.
The Normal Wald test provides one approach, but its normal approximation
can give unreliable results when samples are small or many combinations
of categories are rarely observed.

This thesis develops an Expanded Welch--Satterthwaite test that uses the
same estimated MI difference and standard error as Normal Wald, but
compares the resulting statistic with a Student \(t\) distribution.
Choosing the degrees of freedom requires an estimate of how much the MI
variance estimate changes between samples. The thesis derives a first-order
approximation to this variability from the effect of changes in the
observed category counts.

Simulations compare the tests across sample sizes, table sizes and category
frequencies. They examine how often each test reports a difference when
population MI is equal, and how often it detects a true difference.
An online-shopping dataset illustrates their use.

Expanded Welch reduces excessive false positives in some small-sample
settings, but also rejects less often where Wald already rejects too
rarely. It does not correct bias in the estimated MI difference or errors
in the estimated standard error. The findings therefore do not support
Expanded Welch as a general replacement for Normal Wald.
